\CWHeader{Лабораторная работа \textnumero 2}

\CWProblem{
{\bfseries Вариант структуры данных:} АВЛ-дерево (AVL-tree).

Необходимо создать программную библиотеку, реализующую указанную структуру данных, на основе которой разработать программу-словарь. В словаре каждому ключу, представляющему из себя регистронезависимую последовательность букв английского алфавита длиной не более 256 символов, поставлен в соответствие некоторый номер, от 0 до $2^{64}-1$. Разным словам может быть поставлен в соответствие один и тот же номер.

Программа должна обрабатывать строки входного файла до его окончания. Каждая строка может иметь следующий формат:

\begin{itemize}
    \item \texttt{+ word 34} — добавить слово «word» с номером 34 в словарь. Программа должна вывести строку «OK», если операция прошла успешно, «Exist», если слово уже находится в словаре.
    \item \texttt{- word} — удалить слово «word» из словаря. Программа должна вывести «OK», если слово существовало и было удалено, «NoSuchWord», если слово в словаре не было найдено.
    \item \texttt{word} — найти в словаре слово «word». Программа должна вывести «OK: 34», если слово было найдено. В случае, если слово в словаре не было обнаружено, нужно вывести строку «NoSuchWord».
    \item \texttt{! Save /path/to/file} — сохранить словарь в бинарном компактном представлении на диск. Вывести «OK» или описание ошибки.
    \item \texttt{! Load /path/to/file} — загрузить словарь из файла. В случае успеха вывести «OK», а загруженный словарь должен заменить текущий; в случае неуспеха вывести диагностику, а рабочий словарь должен остаться без изменений.
\end{itemize}

Для всех операций, в случае возникновения системной ошибки, программа должна вывести строку, начинающуюся с «ERROR:» и описывающую на английском языке возникшую ошибку.
}
\pagebreak